\section{Free Lie algebras and universal enveloping Hopf algebras}\label{sec:pbw}
This appendix proves the universal-enveloping statement used by the source article and explains why a calculation in free associative words is a universal Lie identity.

\subsection{The universal algebra and its ordered normal forms}
For a Lie algebra $\gfrak$ over a field of characteristic zero, define
\begin{equation}\label{eq:U-definition}
 U(\gfrak)=T(\gfrak)\big/\bigl\langle x\otimes y-y\otimes x-[x,y]:x,y\in\gfrak\bigr\rangle.
\end{equation}
A linear map $\gfrak\to A$ to an associative algebra extends uniquely to $T(\gfrak)$ by multiplying images of tensors. It factors through \cref{eq:U-definition} exactly when it preserves the Lie bracket, with $A$ given its commutator bracket. This proves the universal property directly. It remains important that the map $\gfrak\to U(\gfrak)$ be injective.

\begin{theorem}[Ordered-monomial theorem]\label{thm:PBW}
Choose a totally ordered basis $(e_i)$ of $\gfrak$. The monomials
$e_{i_1}\cdots e_{i_n}$ with $i_1\leq\cdots\leq i_n$, together with $1$, form a vector-space basis of $U(\gfrak)$. In particular $\gfrak\to U(\gfrak)$ is injective.
\end{theorem}
\begin{proof}
Use the linear rewriting rule
\begin{equation}\label{eq:PBW-rule}
 e_je_i\longmapsto e_ie_j+[e_j,e_i]\qquad(j>i),
\end{equation}
expanding the bracket as a finite linear combination of basis elements. Rewriting terminates: a swapped adjacent inverted pair reduces the inversion number while preserving word length, and a bracket term reduces length. Order these two nonnegative integers lexicographically, with length first; after every replacement each resulting monomial has smaller measure. The multiset extension of this order gives termination for finite linear combinations.

Disjoint replacements commute. The only overlapping ambiguity comes from a word $abc$ with basis indices satisfying $a>b>c$. Reducing the first inverted pair first gives, before normalizing the shorter terms,
\begin{equation*}
 cba+[b,c]a+b[a,c]+[a,b]c.
\end{equation*}
Reducing the second pair first gives
\begin{equation*}
 cba+c[a,b]+[a,c]b+a[b,c].
\end{equation*}
Their difference reduces, using rules on shorter products, to
\begin{align*}
 [[b,c],a]+[b,[a,c]]+[[a,b],c]
 &= [[b,c],a]+[[c,a],b]+[[a,b],c]\
 &=0
\end{align*}
by Jacobi. Thus every possible first conflict can be resolved. Induction on the terminating measure now proves uniqueness of the reduced form: compare the first two reductions, resolve a disjoint or overlapping conflict as above, and use the induction hypothesis on all terms below the initial measure. This is the usual confluence argument written here for this specific system.

Hence every tensor has a unique normal form in the span of ordered words. Each relation in \cref{eq:U-definition}, also when multiplied on either side by words, has zero normal form, because its two sides are a permitted local rewrite. Therefore the whole defining ideal has zero normal form. Conversely, reduction changes a tensor only by elements of that ideal. The quotient is consequently identified with the vector space of ordered normal forms, proving both spanning and linear independence. In degree one the basis elements remain distinct, proving injectivity.
\end{proof}
The proof uses the basic bracket axioms, not a prior matrix representation of $\gfrak$.

\subsection{The free Lie algebra and the free associative algebra}
Let $F(X,Y)$ be the abstract free Lie algebra, constructed from bracketed words by imposing bilinearity, skew symmetry, and Jacobi. Giving a Lie homomorphism from $F(X,Y)$ is exactly the same as choosing two elements in the target Lie algebra: substitution defines the map, and the imposed relations are precisely what makes it well-defined.

The universal property of $U(F(X,Y))$ therefore says that an associative homomorphism from it to an arbitrary associative algebra $A$ is exactly a choice of two elements of $A$. The free associative algebra $\kk\langle X,Y\rangle$ has the same property. More explicitly, sending its two generators to the corresponding elements of $U(F(X,Y))$ defines one map; the free Lie map $F(X,Y)\to\kk\langle X,Y\rangle$ followed by the universal property of $U$ defines the other. Both compositions fix the generators, and hence are identities. Thus
\begin{equation}\label{eq:U-free}
 U(F(X,Y))\cong\kk\langle X,Y\rangle.
\end{equation}
By \cref{thm:PBW}, $F(X,Y)$ injects into this algebra, and its image is exactly the Lie subalgebra generated by $X,Y$. This identifies the concrete $\LL(X,Y)$ used throughout the paper with the abstract free Lie algebra. In particular, an equality of its bracket polynomials verified as an associative-word equality is an identity valid under substitution in \emph{every} Lie algebra.

\subsection{The enveloping Hopf structure}
Define on $\gfrak$ the maps
\begin{equation}\label{eq:U-Hopf}
 \Delta x=x\otimes1+1\otimes x,\qquad
 \eps(x)=0,\qquad S(x)=-x.
\end{equation}
Extend $\Delta,\eps$ as algebra homomorphisms and $S$ as an antihomomorphism. These definitions respect the relations in \cref{eq:U-definition}: for $\Delta$, the mixed tensor-factor terms in $[\Delta x,\Delta y]$ cancel, leaving $\Delta[x,y]$; for $S$, applying it to $xy-yx-[x,y]$ gives the negative of the same relation. Coassociativity and the counit identities hold on generators and hence on all of $U(\gfrak)$. The antipode identities follow from them on generators and induction on products. Thus $U(\gfrak)$ is a Hopf algebra.

Under \cref{eq:U-free}, this is exactly the unshuffle Hopf algebra of \cref{sec:hopf}; degree completion gives its formal noncommutative power-series version. For a general $\gfrak$, ordinary tensor length need not be a grading of $U(\gfrak)$, since a bracket relation changes length. Formal exponentials in that setting can instead be taken in $U(\gfrak)[[t]]$ with $tX,tY$ as arguments. This avoids assuming an inappropriate grading while preserving every coefficient identity.
